% !TEX root = ../../hw05.tex
\begin{proof}[Solution]

\textbf{Claim.} Let $I\subseteq \mathbb{R}$ be an interval with finite endpoints $a$, $b$. If
$f: I \rightarrow \mathbb{R}$ is uniformly continuous then $f$ is bounded on $I$.

Suppose $f$ is uniformly continuous on $I$. By definition, for every
$\varepsilon > 0$ there is a $\delta > 0$ such that
$$
\forall x,y \in I,\quad |x - y| < \delta \implies |f(x) - f(y)| < \varepsilon.
$$
Choose $\varepsilon = 1$. Then there exists a $\delta > 0$ such that
$$
|x - y| < \delta \implies |f(x) - f(y)| < 1.
$$
Our goal is to show that there exists a finite constant $M > 0$ such that
$$
|f(x)| \le M,\quad \forall x \in I.
$$

If $I$ is empty, the claim is immediate. Otherwise, fix $x_0\in I$ once
and for all.

Since $a$ and $b$ are finite endpoints, we consider
an arbitrary point $x\in I$.
We'll take $N$ steps between $x_0$ and $x$ such that $|x_k - x_{k- 1}| < \delta$
and $|f(x_k) - f(x_{k-1})| < \varepsilon$.

% a figure showing this interval
\autoref{fig:hw05:uniform-continuity-chain} illustrates
the finite chain of small steps.

\begin{marginfigure}
  \centering
  % !TEX root = ../../problems/hw05.tex
% Schematic finite chain; the open-interval case is shown.
% Shared physical dimensions for interval and number-line diagrams.
\tikzset{
  interval diagram/.style={
    x=1cm,y=6mm,font=\footnotesize,
    color=black,line width=0.5pt,line cap=round,line join=round,
    >={Stealth[length=4pt,width=3pt]}
  },
  interval endpoint/.style={
    circle,draw=black,line width=0.5pt,
    minimum size=4pt,inner sep=0pt,outer sep=0pt
  },
  interval open/.style={interval endpoint,fill=white},
  interval closed/.style={interval endpoint,fill=black},
  interval label/.style={inner sep=1pt}
}

\begin{tikzpicture}[interval diagram,x=7mm,y=1mm]
  \draw[gray] (0,0) -- (6,0);
  \node[interval open] at (0,0) {};
  \node[interval open] at (6,0) {};
  \node[above=5pt] at (0,0) {$a$};
  \node[above=5pt] at (6,0) {$b$};

  \draw (0.7,0) -- (5.3,0);
  \node[fill=white,inner sep=2pt] at (3.5,0) {$\cdots$};
  \foreach \position in {0.7,1.7,2.7,4.3,5.3}{
    \fill (\position,0) circle[radius=1.6pt];
  }
  \node[below=5pt] at (0.7,0) {$x_0$};
  \node[below=5pt] at (1.7,0) {$x_1$};
  \node[below=5pt] at (2.7,0) {$x_2$};
  \node[above=5pt] at (4.3,0) {$x_{N-1}$};
  \node[below=5pt] at (5.3,0) {$x_N=x$};

  \draw[<->,gray] (0.7,10) -- (1.7,10)
    node[midway,above=4pt,text=black] {$<\delta$};
  \draw[<->,gray] (4.3,10) -- (5.3,10)
    node[midway,above=4pt,text=black] {$<\delta$};
\end{tikzpicture}

  \caption[Chaining small steps across a bounded interval.]{%
    A schematic chain inside $I$, with the open-interval
    case shown. Every successive gap is less than $\delta$, so each
    corresponding output change is less than $1$. The dots omit
    intermediate steps in the finite chain.}
  \label{fig:hw05:uniform-continuity-chain}
\end{marginfigure}

Over our interval $[a,b]$, if we take $N$ steps,
we can arrange for each to have size at most
$$
\dfrac{b - a}{N}.
$$
For uniform continuity, it is enough to choose $N$ so that
$$
\dfrac{b - a}{N} < \delta \iff \dfrac{b - a}{\delta} < N.
$$
Since $b - a$ and $\delta$ are both fixed, we can select
an integer $N\ge2$ that
works for every pair of points in $I$.

Such an integer exists by the Archimedean property. Set
\[
  x_k=x_0+\frac{k}{N}(x-x_0),\qquad k=0,\ldots,N.
\]
Each $x_k$ lies between $x_0$ and $x$, so it belongs to $I$ because $I$
is an interval. For $k=1,\ldots,N$,
\[
  |x_k-x_{k-1}|=\frac{|x-x_0|}{N}
  \le\frac{b-a}{N}<\delta.
\]
Uniform continuity therefore gives $|f(x_k)-f(x_{k-1})|<1$ at every
step. This construction also works when $x\le x_0$.


For the first two steps, we have
$$
|f(x_2) - f(x_0)| = |f(x_2) - \underbrace{f(x_1) \mathbin{-} f(x_1)}_{0} - f(x_0)|.
$$
Using the triangle inequality $|A| + |B| \ge |A + B|$,
\begin{align*}
|f(x_2) - f(x_1) + f(x_1) - f(x_0)|
&= |[f(x_2) - f(x_1)] + [f(x_1) - f(x_0)]| ,\\
\mathrel{\Longrightarrow}\quad |f(x_2) - f(x_0)| &\le |f(x_2) - f(x_1)| + |f(x_1) - f(x_0)|.
\end{align*}
We can continue adding central terms like $f(x_1)$ that cancel out to generalize
this to
$$
|f(x_N) - f(x_{0})| \le \sum_{k=1}^N |f(x_k) - f(x_{k-1})|.
$$
Since we chose $\varepsilon=1$,
$$
\sum_{k=1}^N |f(x_k) - f(x_{k-1})| \mathrel{<} \varepsilon N \mathrel{=} N.
$$
Hence we've established that
\begin{align*}
|f(x_N) - f(x_0)| \le \sum_{k=1}^N |f(x_k) - f(x_{k-1})| &< N,\\
\mathrel{\Longrightarrow}\quad |f(x_N) - f(x_0)| & < N.
\end{align*}
In this case $x_N = x$, so we add and subtract a $f(x_0)$ term to arrive
at an expression for $|f(x)|$ in terms of $f(x_0)$,
\begin{align*}
|f(x)| &= |f(x) - f(x_0) + f(x_0)|,\\
|f(x)| &= |[f(x) - f(x_0)] + f(x_0)|.
\end{align*}
We can apply the triangle inequality again,
\begin{align*}
|f(x)| &= |[f(x) - f(x_0)] + f(x_0)|,\\
&\mathrel{\le} |f(x) - f(x_0)| + |f(x_0)|.
\end{align*}
Combining our previous result with the new inequality
and setting $M=N+|f(x_0)|>0$,
\begin{align*}
|f(x)| & \le |f(x) - f(x_0)| + |f(x_0)|,\\
|f(x)| &< N + |f(x_0)|,\\
|f(x)| &< M,\qquad \text{for every }x \in I.
\end{align*}
Both $N$ and $x_0$ are fixed, so $M$ does not depend on $x$.
Since $N$ and $|f(x_0)|$ are both finite, $f$ is bounded on $I$. Therefore we have
proved that if $f$ is uniformly continuous over $I$, then it is bounded.

To prove $f(x) = \frac{1}{1 +x}$ is continuous but not uniformly continuous, we can
first prove the case for regular continuity, then
apply the contrapositive of the claim.

% insert diagram of the function itself with special focus on the the limit point
% at x = -1
\autoref{fig:hw05:unbounded-reciprocal} shows the graph
near the excluded endpoint $-1$.

\begin{marginfigure}
  \centering
  % !TEX root = ../../problems/hw05.tex
% The branch is restricted to (-1,1); neither endpoint is included.
% Shared physical dimensions for interval and number-line diagrams.
\tikzset{
  interval diagram/.style={
    x=1cm,y=6mm,font=\footnotesize,
    color=black,line width=0.5pt,line cap=round,line join=round,
    >={Stealth[length=4pt,width=3pt]}
  },
  interval endpoint/.style={
    circle,draw=black,line width=0.5pt,
    minimum size=4pt,inner sep=0pt,outer sep=0pt
  },
  interval open/.style={interval endpoint,fill=white},
  interval closed/.style={interval endpoint,fill=black},
  interval label/.style={inner sep=1pt}
}

\begin{tikzpicture}[interval diagram,x=15mm,y=8mm]
  \draw[gray,densely dashed] (-1,0) -- (-1,5.9);
  \draw[->] (-1.22,0) -- (1.25,0) node[right=4pt] {$x$};
  \draw[->] (0,-0.08) -- (0,6.1) node[above=4pt] {$y$};
  \foreach \height in {1,2,4}{
    \draw[gray] (-0.04,\height) -- (0.04,\height);
    \node[left=5pt] at (0,\height) {$\height$};
  }
  \foreach \position in {-1,1}{
    \draw[gray] (\position,-0.04) -- (\position,0.04);
    \node[below=7pt] at (\position,0) {$\position$};
  }
  \node[below left=5pt] at (0,0) {$0$};

  \draw[<-,domain=-0.82:1,samples=120] plot (\x,{1/(1+\x)});
  \foreach \position/\height in {0/1,-0.5/2,-0.75/4}{
    \fill (\position,\height) circle[radius=1.6pt];
  }
  \node[interval open] at (1,0.5) {};
\end{tikzpicture}

  \caption[An unbounded continuous function on an open interval.]{%
    The graph of $f(x)=1/(1+x)$ on $(-1,1)$ rises
    without bound as $x$ approaches $-1$ from the right. The dashed line
    is the vertical asymptote. The open point at $(1,1/2)$ is also
    excluded from the graph.}
  \label{fig:hw05:unbounded-reciprocal}
\end{marginfigure}

The numerator $1$ and denominator $1+x$ are continuous,
and $1+x>0$ on $(-1,1)$.
The quotient of two continuous functions is continuous wherever the denominator is
not equal to zero, so for $x \in (-1, 1)$ the function $f(x)$ is continuous.

When we analyze the limit as $x \rightarrow -1^+$,
$$
\lim_{x \rightarrow -1^+} \dfrac{1}{1 + x} = +\infty,
$$

This means that there is no finite upper bound for $f$ on $(-1,1)$.
More explicitly, for any $B>0$, choose an integer $n>\max\{B,1\}$.
Then
\[
  x_n=-1+\frac1n\in(-1,1),\qquad
  f(x_n)=\frac{1}{1+x_n}=n>B.
\]
Thus $f$ is unbounded on $(-1,1)$.

By the contrapositive of the above claim, if $f$ is not bounded, then it cannot be
uniformly continuous over $(-1,1)$.

\end{proof}
